\section{缺点与变体}

\subsection{Transformer 的主要缺点}

即使 Transformer 很强，它仍然有两个广为人知的短板：第一，self-attention 的计算和显存开销对序列长度是二次的；第二，绝对位置编码未必是最好的顺序表示。

\begin{figure}[H]
\centering
\includegraphics[width=0.95\textwidth]{slides-images/slide-057.jpg}
\caption{Transformer 的主要问题：self-attention 的计算随序列长度二次增长\protect\footnotemark}
\end{figure}
\footnotetext{Slides 第58页。}

\begin{importantbox}{两个主要方向}
\begin{enumerate}
    \item 降低 $O(T^2)$ self-attention 的成本
    \item 改进位置表示，让模型更好地利用顺序与结构信息
\end{enumerate}
\end{importantbox}

\subsection{位置表示的后续工作}

讲义里提到的几个方向包括 relative linear position attention、dependency syntax-based position 和 rotary embeddings。它们都试图让位置不只是一个静态的绝对 index，而是更贴近词与词之间的关系。

\begin{figure}[H]
\centering
\includegraphics[width=0.95\textwidth]{slides-images/slide-058.jpg}
\caption{一些改进位置表示的思路：相对位置、句法位置、旋转位置编码等\protect\footnotemark}
\end{figure}
\footnotetext{Slides 第59页。}

\subsection{高效 attention 变体}

为了摆脱 $O(T^2)$，研究者提出了许多高效 Transformer 变体。例如 Linformer 通过把序列长度维度投影到低维空间减少成本，BigBird 则用局部窗口、全局 token 和随机连接替代全连接注意力。

\begin{figure}[H]
\centering
\includegraphics[width=0.95\textwidth]{slides-images/slide-059.jpg}
\caption{Linformer：把序列长度投影到低维表示，降低 attention 成本\protect\footnotemark}
\end{figure}
\footnotetext{Slides 第60页。}

\begin{figure}[H]
\centering
\includegraphics[width=0.95\textwidth]{slides-images/slide-060.jpg}
\caption{BigBird：用局部、全局和随机注意力模式替换全连接 attention\protect\footnotemark}
\end{figure}
\footnotetext{Slides 第61页。}

\begin{warningbox}{变体不一定真的更好}
讲义最后提醒：很多 Transformer 修改并不会带来显著收益。也就是说，改结构并不天然意味着更强，是否 transfer 到实际任务中仍然需要实证验证。
\end{warningbox}

\begin{figure}[H]
\centering
\includegraphics[width=0.95\textwidth]{slides-images/slide-061.jpg}
\caption{总结性结论：很多修改并没有显著提升性能\protect\footnotemark}
\end{figure}
\footnotetext{Slides 第62页。}

\subsection{本章小结}

Transformer 的成功并不意味着它已经没有缺陷。它仍然面对 quadratic cost 和位置表示设计的问题。后续研究主要围绕两条线展开：一条是让 attention 更省算力，另一条是让位置编码更符合语言结构。

%% ========== 总结与延伸 ========== 

